%% kapa: exam=1 round=22 year=2011 period=1 form=A subject=economics subject_ko=경제원론 seq=19 number=59
%% source: Q-Net 감정평가사 기출문제 (cst00309) artlSeq=5036790 file=extracted/1008/5036790/2011년 제22회 감정평가사 1차시험 기출문제/1교시.hwp
%% conversion: kordoc 4.24.51 -> tools/gen1.py (auto); review=proofread-sonnet; flags=none
\begin{question}{59}
A은행의 T-계정은 다음과 같다.

{\small\centering
\begin{tabular}{|l|l|}
\hline
\multicolumn{1}{|c|}{자산} & \multicolumn{1}{c|}{부채} \\
\hline
\makecell[l]{지급준비금 1,000억원\\대출 3,000억원} & 예금 4,000억원 \\
\hline
\end{tabular}\par}

예금에 대한 법정지급준비율이 10\%이고, A은행을 제외한 다른 은행들은 초과 지급준비금을 보유하지 않는다. A은행이 지급준비금을 법정지급준비금 수준까지 줄인다면 최대로 가능한 통화량 증가액은? (단, 민간의 현금보유비율은 0) \par
\begin{choices}
  \item 600억원
  \item 1,000억원
  \item 4,000억원
  \item 6,000억원
  \item 1조원
\end{choices}
\end{question}
\answer{4}
